\documentclass[10pt,a4paper]{article}
\usepackage{fontspec,amsmath,amssymb,graphicx}
\usepackage[margin=18mm]{geometry}
\setmainfont{Times New Roman}
\setlength{\parindent}{0pt}
\setlength{\parskip}{5pt}
\allowdisplaybreaks
\sloppy
\emergencystretch=2em
\begin{document}
\section*{VD\_4\_5\_2\_01 — Lời giải bổ sung}
\subsection*{Kiến thức và trình tự giải}
\textbf{Bước 1 — Loại đạo hàm cấp một bằng tính đạo hàm đầy đủ.} Với $u_t=Ku_{xx}+Au_x+Bu$, đặt $u=e^{-hx}v$, $h=A/(2K)$. Khi đó


\[
u_t=e^{-hx}v_t,\quad u_x=e^{-hx}(v_x-hv),\quad
u_{xx}=e^{-hx}(v_{xx}-2hv_x+h^2v).
\]


Thay vào PDE, hệ số $v_x$ là $A-2Kh=0$, hệ số $v$ là $Kh^2-Ah+B=B-A^2/(4K)$. Vì thế $v_t=Kv_{xx}+Cv$, $C=B-A^2/(4K)$; có thể tiếp tục đặt $v=e^{Ct}z$ để $z_t=Kz_{xx}$.

\textbf{Bước 2 — Đổi dữ liệu và biên cùng phép đổi ẩn.} $v(x,0)=e^{hx}f(x)$. Dirichlet thuần nhất giữ nguyên; Neumann biến thành $v_x-hv=0$, nên không giữ nguyên một cơ sở cosine nếu $A\ne0$. Với Dirichlet:


\[
u=e^{-hx}\sum_{n\ge1}\left(\frac2L\int_0^Le^{hy}f(y)\sin\frac{n\pi y}Ldy\right)
e^{[C-K(n\pi/L)^2]t}\sin\frac{n\pi x}L.
\]


\textbf{Bước 3 — Tính hệ số mũ–đa thức.} Dùng nguyên hàm


\[
\int e^{hy}\sin ky\,dy
=\frac{e^{hy}(h\sin ky-k\cos ky)}{h^2+k^2}.
\]


Với $y e^{hy}$ hoặc $y(1-y)e^{hy}$, lấy đạo hàm theo $h$ của tích phân đã tính trên cận cố định; trình bày đạo hàm để tìm hệ số thực. Cách này tránh để lại một tích phân sơ cấp chưa thực hiện.

\textbf{Bước 4 — Hệ số không gian biến thiên.} Với $x^2u_{xx}$ trên $x>0$, đặt $z=\ln x$: nếu $U(z,t)=u(e^z,t)$ thì $u_x=U_z/x$, $x^2u_{xx}=U_{zz}-U_z$. Đổi miền và biên sang biến $z$ rồi dùng bước 1. Biên không thuần nhất nên tìm trạng thái dừng trước khi khai triển.

\textbf{Bước 5 — Duy nhất và ổn định.} Với hiệu nghiệm Dirichlet, $I'=-K\int w_x^2+2BI\le2BI$ suy $w=0$ khi $I(0)=0$. Giới hạn phụ thuộc các số $C-K\lambda_n$ có hệ số khác không: âm thì giảm, bằng không thì tồn tại mode dừng, dương thì tăng. Phải kiểm tra tính thống nhất giữa miền PDE và vị trí biên trước khi dùng công thức.
\subsection*{Lời giải áp dụng cho bài}
\textbf{Bước 1. Tính phép đổi ẩn.}
Đặt $u=e^{-hx}v$, $h=A/(2K)$. Khi đó $u_x=e^{-hx}(v_x-hv)$, $u_{xx}=e^{-hx}(v_{xx}-2hv_x+h^2v)$, $u_t=e^{-hx}v_t$. Thay vào $u_t=Ku_{xx}+Au_x+Bu$ cho
\[v_t=Kv_{xx}+(B-A^2/(4K))v,\quad v(x,0)=e^{hx}f(x).\]
Hai biên Dirichlet vẫn bằng $0$, nên $u=e^{-hx}\sum c_ne^{[B-A^2/(4K)-Kn^2]t}\sin nx$ khi $L=\pi$.

\textbf{Bước 2. Câu (i), $K=1,A=0,B=-1,f=x$.}
\[c_n=\frac2\pi\int_0^\pi x\sin nx\,dx
=\frac2\pi\left[-\frac{x\cos nx}n+\frac{\sin nx}{n^2}\right]_0^\pi
=\frac{2(-1)^{n+1}}n.\]
Vì vậy $u=\sum_{n\ge1}2(-1)^{n+1}e^{-(n^2+1)t}\sin nx/n$.

\textbf{Bước 3. Câu (ii), $K=1,A=-1,B=0,f=1$.}
$h=-1/2$, nên dữ liệu của $v$ là $e^{-x/2}$. Nguyên hàm $\int e^{hx}\sin nx=e^{hx}(h\sin nx-n\cos nx)/(h^2+n^2)$ cho
\[c_n=\frac2\pi\left[\frac{e^{-x/2}(-\sin nx/2-n\cos nx)}{n^2+1/4}\right]_0^\pi
=\frac{2n(1-(-1)^ne^{-\pi/2})}{\pi(n^2+1/4)}.\]
Suy $u=e^{x/2}\sum c_ne^{-(n^2+1/4)t}\sin nx$.

\textbf{Bước 4. Duy nhất và kiểm tra.}
Với hiệu $w$ có dữ liệu đầu $0$, biên $0$,
\begin{align*}
I'&=K[ww_x]_0^\pi-K\int w_x^2+\frac A2[w^2]_0^\pi+B\int w^2\\
&=-K\int w_x^2+2BI\le0
\end{align*}
trong cả hai câu vì $B\le0$. $I(0)=0$ suy $w=0$. Phép đổi ẩn và các ODE mode kiểm tra PDE; sine kiểm tra biên; Fourier trả lại dữ liệu trong miền. $f(\pi)=\pi$ ở (i), và $f(0)=f(\pi)=1$ ở (ii) không khớp biên $0$, nên không liên tục tại các góc này. Hai nghiệm trơn khi $t>0$ và đều tiến về $0$.\end{document}